% =====================================================================
%  Artifact Appendix for the TRAIL paper.
%
%  Drop-in LaTeX. Place near the end of the paper, after \bibliography, e.g.
%
%      \appendix
%      % =====================================================================
%  Artifact Appendix for the TRAIL paper.
%
%  Drop-in LaTeX. Place near the end of the paper, after \bibliography, e.g.
%
%      \appendix
%      % =====================================================================
%  Artifact Appendix for the TRAIL paper.
%
%  Drop-in LaTeX. Place near the end of the paper, after \bibliography, e.g.
%
%      \appendix
%      % =====================================================================
%  Artifact Appendix for the TRAIL paper.
%
%  Drop-in LaTeX. Place near the end of the paper, after \bibliography, e.g.
%
%      \appendix
%      \input{experiments/ae/artifact_appendix}
%
%  It uses only \section/\subsection, itemize, \texttt, and \url (hyperref).
%  If your template already issues \appendix, delete the \appendix line below.
%
%  FILL-INS (search for <<>>): paper title and, if you mint one, a Zenodo DOI.
% =====================================================================

\appendix
\section{Artifact Appendix}

\subsection{Abstract}

This artifact reproduces the main results of TRAIL, a temporal TLB
prefetcher that makes the page table self-prefetching: each translation
carries its likely successors, so prefetching a translation simultaneously
delivers the next prediction. The artifact evaluates TRAIL against prior
TLB prefetchers (ASP, Stride/NextPage, DP, Recency, ATP, Berti) and the
Perfect-L2TLB upper bound on the Sniper architectural simulator extended
with the MimicOS operating-system model. It reproduces the single-core
head-to-head results at an 8\,MB and a 2\,MB last-level-cache NUCA
(Figures~12 and~13), the in-PTE and side-car payload-budget sweeps
(Tables~5 and~6), the L2-TLB prefetch-queue sensitivity study
(Figure~20), and the four-core multi-programmed study (Figure~22). It also
reproduces the temporal-locality motivation figures (Figures~4, 5, 6, 8,
and~9) directly from page-table-walk dumps. Evaluation is trace-driven, so
no proprietary tools (Intel Pin/SDE) are required: a single script installs
dependencies, builds the simulator, downloads the public traces, and runs a
sanity check; a three-command launch/watch/results workflow then reproduces
each claim and renders the corresponding figure or table.

\subsection{Artifact check-list (meta-information)}

\begin{itemize}
  \item \textbf{Program:} Sniper multi-core simulator with the MimicOS OS
        model and the TRAIL prefetcher (C++17); Python analysis and plotting
        scripts. All included.
  \item \textbf{Compilation:} \texttt{g++} (C++17); the instruction decoder
        (libxed) is bundled and built automatically. No Intel Pin/SDE or
        libtorch is needed for the trace-replay build.
  \item \textbf{Run-time environment:} Linux (tested on Ubuntu 20.04 and
        22.04). Root/\texttt{sudo} only to install system packages.
  \item \textbf{Hardware:} any x86-64 machine. A SLURM cluster is recommended
        for the full run but not required; every step also runs locally.
  \item \textbf{Metrics:} instructions-per-cycle (IPC) and its speedup over a
        no-prefetch baseline. Single-core speedups are geometric means over the
        workload suite; the multi-core metric is an equal-work harmonic mean
        across the four cores.
  \item \textbf{Output:} per-claim markdown tables and rendered figures/tables
        (PDF\,+\,PNG) under \texttt{experiments/ae/ae\_out/}.
  \item \textbf{Experiments:} six simulation claims (\texttt{head8mb},
        \texttt{head2mb}, \texttt{table5}, \texttt{table6}, \texttt{pqsweep},
        \texttt{multicore}) plus the trace-free motivation analysis, driven by
        the provided bash harness.
  \item \textbf{How much disk space required (approximately)?} $\sim$300\,GB
        ($\sim$250\,GB of traces plus the build and results).
  \item \textbf{How much time is needed to prepare workflow (approximately)?}
        $\sim$30--60\,min: a 2--3\,min build plus the (resumable) trace
        download, which dominates and depends on bandwidth.
  \item \textbf{How much time is needed to complete experiments
        (approximately)?} The full set is $\approx$21{,}000 short simulations.
        On a $\sim$1300-core cluster the entire set finishes within
        $\sim$1\,day; the longest single claim (\texttt{table6}) takes
        $\sim$10\,h. A single claim, e.g.\ the headline \texttt{head8mb}, is
        far cheaper; an \texttt{--icount} flag shrinks every job to a
        minutes-long smoke test.
  \item \textbf{Publicly available?} Yes. Source on GitHub; traces and
        page-table-walk dumps as public Hugging Face datasets (no token
        required).
  \item \textbf{Code licenses (if publicly available)?} See the repository
        \texttt{LICENSE}. \emph{<<confirm license, e.g.\ MIT>>}
  \item \textbf{Data licenses (if publicly available)?} Public Hugging Face
        datasets (traces and dumps).
  \item \textbf{Workflow automation framework used?} A self-contained bash
        harness (\texttt{reproduce.sh} for setup, then
        \texttt{ae\_launch.sh}\,/\,\texttt{ae\_watch.sh}\,/\,%
        \texttt{ae\_results.sh}, or \texttt{ae\_run\_all.sh} to run every claim
        in parallel).
  \item \textbf{Archived (provide DOI)?} \emph{<<add Zenodo DOI if minted;
        otherwise the GitHub tag + Hugging Face datasets>>}
\end{itemize}

\subsection{Description}

\subsubsection{How to access}
Clone the artifact branch of the repository:
\begin{quote}\small
\texttt{git clone -{}-branch trail-artifact-release \textbackslash}\\
\texttt{\ \ https://github.com/CMU-SAFARI/Virtuoso.git}
\end{quote}
The traces and page-table-walk dumps are separate public Hugging Face
datasets and are fetched by the setup script:
\begin{itemize}
  \item Traces + trace-lists: \url{https://huggingface.co/datasets/konkanello/trail_traces}
  \item Motivation dumps: \url{https://huggingface.co/datasets/konkanello/trail_ptw_dumps}
\end{itemize}
All artifact material lives under \texttt{experiments/ae/}.

\subsubsection{Hardware dependencies}
An x86-64 machine. The build and the sanity check need $\sim$8\,GB of RAM;
each simulation uses a few GB, so a machine (or cluster) with more memory and
cores runs more jobs in parallel. $\sim$300\,GB of free disk is required for
the traces, build, and results. A SLURM cluster is optional but strongly
recommended to complete the full set in reasonable wall-clock time.

\subsubsection{Software dependencies}
Linux with a C++17 toolchain and Python~3.8+. All system packages, plus the
Hugging Face CLI (trace download) and matplotlib (figures), are installed by
\texttt{experiments/ae/lib/install\_deps.sh} on Debian/Ubuntu (equivalents
exist on other distributions). Network access is needed to download the traces
and to fetch the decoder during the build. No license-gated or proprietary
tools are required.

\subsubsection{Benchmarks / data sets}
The workload suite spans server, graph-analytics, machine-learning, HPC, and
industrial traces. The public trace dataset ships the trace files together
with the trace-lists; the setup script wires them into the repository. The
motivation figures use a separate $\sim$3\,GB dataset of per-workload
page-table-walk dumps.

\subsection{Installation}

\begin{enumerate}
  \item Install dependencies:
\begin{quote}\small
\texttt{bash experiments/ae/lib/install\_deps.sh}
\end{quote}
  \item Build, download traces, and sanity-check (resumable; re-run to resume
        after any interruption):
\begin{quote}\small
\texttt{bash experiments/ae/reproduce.sh -{}-skip-deps}
\end{quote}
        This performs a trace-replay build of the simulator
        (\texttt{make SNIPER\_TRACE\_ONLY=1 replay}; no Pin/SDE/libtorch),
        downloads the public trace dataset, resolves the trace-lists into
        \texttt{experiments/vm\_tlist/}, and runs a few short simulations on
        randomly chosen traces, printing \texttt{[PASS] <trace> IPC=...} for
        each. It then stops and prints the commands for Step~3. This step alone
        confirms that the artifact builds and runs end-to-end.
\end{enumerate}

\subsection{Experiment workflow}

Each claim maps to one paper figure or table. Results and rendered figures are
written to \texttt{experiments/ae/ae\_out/}.

\emph{All claims in parallel (recommended on a cluster):}
\begin{quote}\small
\texttt{bash experiments/ae/ae\_run\_all.sh -{}-mode slurm \textbackslash}\\
\texttt{\ \ [-{}-partitions <p>]}\qquad\textit{\# launch every claim}\\
\texttt{bash experiments/ae/ae\_run\_all.sh -{}-status}\quad\textit{\# progress}\\
\texttt{bash experiments/ae/ae\_run\_all.sh -{}-results}\quad\textit{\# figures/tables}
\end{quote}
\texttt{-{}-partitions} is optional (omit it to use the cluster default; nothing
is then passed to \texttt{sbatch}).

\emph{One claim at a time (launch $\rightarrow$ watch $\rightarrow$ results):}
\begin{quote}\small
\texttt{bash experiments/ae/ae\_launch.sh\ -{}-claim head8mb -{}-mode slurm}\\
\texttt{bash experiments/ae/ae\_watch.sh\ \ -{}-claim head8mb}\\
\texttt{bash experiments/ae/ae\_results.sh -{}-claim head8mb}\quad\textit{\# after .DONE}
\end{quote}
The watcher self-detaches and maintains a human-readable
\texttt{ae\_out/<claim>.status}; when every job is accounted for it writes a
pass/fail report to \texttt{ae\_out/<claim>.DONE}, and \texttt{ae\_results.sh}
then parses the results and renders the figure/table. Launching is
validity-based and resumable: only jobs without a valid \texttt{sim.stats} are
(re)submitted, so an interrupted run resumes cleanly. On a single machine, use
\texttt{-{}-mode local} one claim at a time; \texttt{-{}-icount N} shrinks every
job's instruction budget for a quick end-to-end test.

The claim-to-paper mapping is:

\begin{center}\small
\begin{tabular}{@{}llr@{}}
\toprule
\texttt{-{}-claim} & reproduces & jobs \\
\midrule
\texttt{head8mb}   & Fig.~12 (bottom, 8\,MB) + Fig.~13 & 2761 \\
\texttt{head2mb}   & Fig.~12 (top, 2\,MB)              & 2761 \\
\texttt{table5}    & Table~5 (in-PTE payload sweep)    & 5271 \\
\texttt{table6}    & Table~6 (side-car payload sweep)  & 6275 \\
\texttt{pqsweep}   & Fig.~20 (PQ-size sensitivity)     & 3012 \\
\texttt{multicore} & Fig.~22 (4-core, 100 mixes)       & 1000 \\
\midrule
\textbf{all}       &                                   & \textbf{21080} \\
\bottomrule
\end{tabular}
\end{center}

\subsection{Evaluation and expected results}

Each claim writes a markdown table (\texttt{ae\_out/<claim>.md}) alongside the
rendered figure/table. Exact numbers vary slightly with the machine, but every
claim should closely match the paper, and in all cases \textbf{TRAIL should
beat every prior prefetcher} and move toward the Perfect-L2TLB upper bound.
Speedups are geometric means over the suite (multi-core is an equal-work
harmonic mean across the four cores):

\begin{center}\small
\begin{tabular}{@{}ll@{}}
\toprule
claim & expected (approx.) \\
\midrule
\texttt{head8mb} (Fig.~12, 8\,MB) & TRAIL $\approx +4.7\%$ over no-prefetch \\
                                  & ($+2.4\%$ over ASP); Perfect $\approx +11\%$ \\
\texttt{head2mb} (Fig.~12, 2\,MB) & TRAIL $\approx +5.1\%$ ($+2.8\%$ over ASP); \\
                                  & Perfect $\approx +16.5\%$ \\
\texttt{table5}  (Table~5)        & in-PTE TRAIL grows with the budget, \\
                                  & up to $\approx +2.4\%$ over ASP \\
\texttt{table6}  (Table~6)        & side-car TRAIL grows with the budget, \\
                                  & up to $\approx +2.5\%$ over ASP \\
\texttt{pqsweep} (Fig.~20)        & TRAIL $\approx +2.2$--$2.4\%$ over same- \\
                                  & size ASP across every PQ size (64--1024) \\
\texttt{multicore} (Fig.~22)      & TRAIL $\approx +11\%$ harmonic mean (best \\
                                  & prior $\approx +5\%$); Perfect $\approx +23\%$ \\
\bottomrule
\end{tabular}
\end{center}

\subsection{Motivation figures (Figures 4, 5, 6, 8, 9)}

The temporal-locality motivation figures are produced from the page-table-walk
dumps, without running any simulation. Download the dumps and run the analysis
(SLURM or local, using up to \texttt{\$(nproc)} cores):
\begin{quote}\small
\texttt{hf download konkanello/trail\_ptw\_dumps \textbackslash}\\
\texttt{\ \ -{}-repo-type dataset -{}-local-dir ./ptw\_bundle}\\[2pt]
\texttt{bash experiments/ae/motivation/run\_motivation.sh \textbackslash}\\
\texttt{\ \ -{}-dumps ./ptw\_bundle -{}-mode local -{}-jobs \$(nproc)}
\end{quote}
This writes \texttt{figure4}--\texttt{figure9} (PDF\,+\,PNG) under
\texttt{experiments/ae/motivation/motivation\_out/}. The analysis is resumable:
one JSON per workload, re-plotted on any re-run.

\subsection{Notes}

\begin{itemize}
  \item Every phase is resumable. Re-running \texttt{reproduce.sh}, a launch, or
        the motivation driver skips whatever is already complete.
  \item Reviewers on a time budget should start with the setup sanity check
        (Installation, Step~2) and the headline \texttt{head8mb} claim, which is
        the cheapest full claim; \texttt{-{}-icount} makes any claim a
        minutes-long end-to-end smoke test.
  \item The full artifact README is at \texttt{experiments/ae/README.md}, and
        the motivation README at \texttt{experiments/ae/motivation/README.md}.
\end{itemize}

%
%  It uses only \section/\subsection, itemize, \texttt, and \url (hyperref).
%  If your template already issues \appendix, delete the \appendix line below.
%
%  FILL-INS (search for <<>>): paper title and, if you mint one, a Zenodo DOI.
% =====================================================================

\appendix
\section{Artifact Appendix}

\subsection{Abstract}

This artifact reproduces the main results of TRAIL, a temporal TLB
prefetcher that makes the page table self-prefetching: each translation
carries its likely successors, so prefetching a translation simultaneously
delivers the next prediction. The artifact evaluates TRAIL against prior
TLB prefetchers (ASP, Stride/NextPage, DP, Recency, ATP, Berti) and the
Perfect-L2TLB upper bound on the Sniper architectural simulator extended
with the MimicOS operating-system model. It reproduces the single-core
head-to-head results at an 8\,MB and a 2\,MB last-level-cache NUCA
(Figures~12 and~13), the in-PTE and side-car payload-budget sweeps
(Tables~5 and~6), the L2-TLB prefetch-queue sensitivity study
(Figure~20), and the four-core multi-programmed study (Figure~22). It also
reproduces the temporal-locality motivation figures (Figures~4, 5, 6, 8,
and~9) directly from page-table-walk dumps. Evaluation is trace-driven, so
no proprietary tools (Intel Pin/SDE) are required: a single script installs
dependencies, builds the simulator, downloads the public traces, and runs a
sanity check; a three-command launch/watch/results workflow then reproduces
each claim and renders the corresponding figure or table.

\subsection{Artifact check-list (meta-information)}

\begin{itemize}
  \item \textbf{Program:} Sniper multi-core simulator with the MimicOS OS
        model and the TRAIL prefetcher (C++17); Python analysis and plotting
        scripts. All included.
  \item \textbf{Compilation:} \texttt{g++} (C++17); the instruction decoder
        (libxed) is bundled and built automatically. No Intel Pin/SDE or
        libtorch is needed for the trace-replay build.
  \item \textbf{Run-time environment:} Linux (tested on Ubuntu 20.04 and
        22.04). Root/\texttt{sudo} only to install system packages.
  \item \textbf{Hardware:} any x86-64 machine. A SLURM cluster is recommended
        for the full run but not required; every step also runs locally.
  \item \textbf{Metrics:} instructions-per-cycle (IPC) and its speedup over a
        no-prefetch baseline. Single-core speedups are geometric means over the
        workload suite; the multi-core metric is an equal-work harmonic mean
        across the four cores.
  \item \textbf{Output:} per-claim markdown tables and rendered figures/tables
        (PDF\,+\,PNG) under \texttt{experiments/ae/ae\_out/}.
  \item \textbf{Experiments:} six simulation claims (\texttt{head8mb},
        \texttt{head2mb}, \texttt{table5}, \texttt{table6}, \texttt{pqsweep},
        \texttt{multicore}) plus the trace-free motivation analysis, driven by
        the provided bash harness.
  \item \textbf{How much disk space required (approximately)?} $\sim$300\,GB
        ($\sim$250\,GB of traces plus the build and results).
  \item \textbf{How much time is needed to prepare workflow (approximately)?}
        $\sim$30--60\,min: a 2--3\,min build plus the (resumable) trace
        download, which dominates and depends on bandwidth.
  \item \textbf{How much time is needed to complete experiments
        (approximately)?} The full set is $\approx$21{,}000 short simulations.
        On a $\sim$1300-core cluster the entire set finishes within
        $\sim$1\,day; the longest single claim (\texttt{table6}) takes
        $\sim$10\,h. A single claim, e.g.\ the headline \texttt{head8mb}, is
        far cheaper; an \texttt{--icount} flag shrinks every job to a
        minutes-long smoke test.
  \item \textbf{Publicly available?} Yes. Source on GitHub; traces and
        page-table-walk dumps as public Hugging Face datasets (no token
        required).
  \item \textbf{Code licenses (if publicly available)?} See the repository
        \texttt{LICENSE}. \emph{<<confirm license, e.g.\ MIT>>}
  \item \textbf{Data licenses (if publicly available)?} Public Hugging Face
        datasets (traces and dumps).
  \item \textbf{Workflow automation framework used?} A self-contained bash
        harness (\texttt{reproduce.sh} for setup, then
        \texttt{ae\_launch.sh}\,/\,\texttt{ae\_watch.sh}\,/\,%
        \texttt{ae\_results.sh}, or \texttt{ae\_run\_all.sh} to run every claim
        in parallel).
  \item \textbf{Archived (provide DOI)?} \emph{<<add Zenodo DOI if minted;
        otherwise the GitHub tag + Hugging Face datasets>>}
\end{itemize}

\subsection{Description}

\subsubsection{How to access}
Clone the artifact branch of the repository:
\begin{quote}\small
\texttt{git clone -{}-branch trail-artifact-release \textbackslash}\\
\texttt{\ \ https://github.com/CMU-SAFARI/Virtuoso.git}
\end{quote}
The traces and page-table-walk dumps are separate public Hugging Face
datasets and are fetched by the setup script:
\begin{itemize}
  \item Traces + trace-lists: \url{https://huggingface.co/datasets/konkanello/trail_traces}
  \item Motivation dumps: \url{https://huggingface.co/datasets/konkanello/trail_ptw_dumps}
\end{itemize}
All artifact material lives under \texttt{experiments/ae/}.

\subsubsection{Hardware dependencies}
An x86-64 machine. The build and the sanity check need $\sim$8\,GB of RAM;
each simulation uses a few GB, so a machine (or cluster) with more memory and
cores runs more jobs in parallel. $\sim$300\,GB of free disk is required for
the traces, build, and results. A SLURM cluster is optional but strongly
recommended to complete the full set in reasonable wall-clock time.

\subsubsection{Software dependencies}
Linux with a C++17 toolchain and Python~3.8+. All system packages, plus the
Hugging Face CLI (trace download) and matplotlib (figures), are installed by
\texttt{experiments/ae/lib/install\_deps.sh} on Debian/Ubuntu (equivalents
exist on other distributions). Network access is needed to download the traces
and to fetch the decoder during the build. No license-gated or proprietary
tools are required.

\subsubsection{Benchmarks / data sets}
The workload suite spans server, graph-analytics, machine-learning, HPC, and
industrial traces. The public trace dataset ships the trace files together
with the trace-lists; the setup script wires them into the repository. The
motivation figures use a separate $\sim$3\,GB dataset of per-workload
page-table-walk dumps.

\subsection{Installation}

\begin{enumerate}
  \item Install dependencies:
\begin{quote}\small
\texttt{bash experiments/ae/lib/install\_deps.sh}
\end{quote}
  \item Build, download traces, and sanity-check (resumable; re-run to resume
        after any interruption):
\begin{quote}\small
\texttt{bash experiments/ae/reproduce.sh -{}-skip-deps}
\end{quote}
        This performs a trace-replay build of the simulator
        (\texttt{make SNIPER\_TRACE\_ONLY=1 replay}; no Pin/SDE/libtorch),
        downloads the public trace dataset, resolves the trace-lists into
        \texttt{experiments/vm\_tlist/}, and runs a few short simulations on
        randomly chosen traces, printing \texttt{[PASS] <trace> IPC=...} for
        each. It then stops and prints the commands for Step~3. This step alone
        confirms that the artifact builds and runs end-to-end.
\end{enumerate}

\subsection{Experiment workflow}

Each claim maps to one paper figure or table. Results and rendered figures are
written to \texttt{experiments/ae/ae\_out/}.

\emph{All claims in parallel (recommended on a cluster):}
\begin{quote}\small
\texttt{bash experiments/ae/ae\_run\_all.sh -{}-mode slurm \textbackslash}\\
\texttt{\ \ [-{}-partitions <p>]}\qquad\textit{\# launch every claim}\\
\texttt{bash experiments/ae/ae\_run\_all.sh -{}-status}\quad\textit{\# progress}\\
\texttt{bash experiments/ae/ae\_run\_all.sh -{}-results}\quad\textit{\# figures/tables}
\end{quote}
\texttt{-{}-partitions} is optional (omit it to use the cluster default; nothing
is then passed to \texttt{sbatch}).

\emph{One claim at a time (launch $\rightarrow$ watch $\rightarrow$ results):}
\begin{quote}\small
\texttt{bash experiments/ae/ae\_launch.sh\ -{}-claim head8mb -{}-mode slurm}\\
\texttt{bash experiments/ae/ae\_watch.sh\ \ -{}-claim head8mb}\\
\texttt{bash experiments/ae/ae\_results.sh -{}-claim head8mb}\quad\textit{\# after .DONE}
\end{quote}
The watcher self-detaches and maintains a human-readable
\texttt{ae\_out/<claim>.status}; when every job is accounted for it writes a
pass/fail report to \texttt{ae\_out/<claim>.DONE}, and \texttt{ae\_results.sh}
then parses the results and renders the figure/table. Launching is
validity-based and resumable: only jobs without a valid \texttt{sim.stats} are
(re)submitted, so an interrupted run resumes cleanly. On a single machine, use
\texttt{-{}-mode local} one claim at a time; \texttt{-{}-icount N} shrinks every
job's instruction budget for a quick end-to-end test.

The claim-to-paper mapping is:

\begin{center}\small
\begin{tabular}{@{}llr@{}}
\toprule
\texttt{-{}-claim} & reproduces & jobs \\
\midrule
\texttt{head8mb}   & Fig.~12 (bottom, 8\,MB) + Fig.~13 & 2761 \\
\texttt{head2mb}   & Fig.~12 (top, 2\,MB)              & 2761 \\
\texttt{table5}    & Table~5 (in-PTE payload sweep)    & 5271 \\
\texttt{table6}    & Table~6 (side-car payload sweep)  & 6275 \\
\texttt{pqsweep}   & Fig.~20 (PQ-size sensitivity)     & 3012 \\
\texttt{multicore} & Fig.~22 (4-core, 100 mixes)       & 1000 \\
\midrule
\textbf{all}       &                                   & \textbf{21080} \\
\bottomrule
\end{tabular}
\end{center}

\subsection{Evaluation and expected results}

Each claim writes a markdown table (\texttt{ae\_out/<claim>.md}) alongside the
rendered figure/table. Exact numbers vary slightly with the machine, but every
claim should closely match the paper, and in all cases \textbf{TRAIL should
beat every prior prefetcher} and move toward the Perfect-L2TLB upper bound.
Speedups are geometric means over the suite (multi-core is an equal-work
harmonic mean across the four cores):

\begin{center}\small
\begin{tabular}{@{}ll@{}}
\toprule
claim & expected (approx.) \\
\midrule
\texttt{head8mb} (Fig.~12, 8\,MB) & TRAIL $\approx +4.7\%$ over no-prefetch \\
                                  & ($+2.4\%$ over ASP); Perfect $\approx +11\%$ \\
\texttt{head2mb} (Fig.~12, 2\,MB) & TRAIL $\approx +5.1\%$ ($+2.8\%$ over ASP); \\
                                  & Perfect $\approx +16.5\%$ \\
\texttt{table5}  (Table~5)        & in-PTE TRAIL grows with the budget, \\
                                  & up to $\approx +2.4\%$ over ASP \\
\texttt{table6}  (Table~6)        & side-car TRAIL grows with the budget, \\
                                  & up to $\approx +2.5\%$ over ASP \\
\texttt{pqsweep} (Fig.~20)        & TRAIL $\approx +2.2$--$2.4\%$ over same- \\
                                  & size ASP across every PQ size (64--1024) \\
\texttt{multicore} (Fig.~22)      & TRAIL $\approx +11\%$ harmonic mean (best \\
                                  & prior $\approx +5\%$); Perfect $\approx +23\%$ \\
\bottomrule
\end{tabular}
\end{center}

\subsection{Motivation figures (Figures 4, 5, 6, 8, 9)}

The temporal-locality motivation figures are produced from the page-table-walk
dumps, without running any simulation. Download the dumps and run the analysis
(SLURM or local, using up to \texttt{\$(nproc)} cores):
\begin{quote}\small
\texttt{hf download konkanello/trail\_ptw\_dumps \textbackslash}\\
\texttt{\ \ -{}-repo-type dataset -{}-local-dir ./ptw\_bundle}\\[2pt]
\texttt{bash experiments/ae/motivation/run\_motivation.sh \textbackslash}\\
\texttt{\ \ -{}-dumps ./ptw\_bundle -{}-mode local -{}-jobs \$(nproc)}
\end{quote}
This writes \texttt{figure4}--\texttt{figure9} (PDF\,+\,PNG) under
\texttt{experiments/ae/motivation/motivation\_out/}. The analysis is resumable:
one JSON per workload, re-plotted on any re-run.

\subsection{Notes}

\begin{itemize}
  \item Every phase is resumable. Re-running \texttt{reproduce.sh}, a launch, or
        the motivation driver skips whatever is already complete.
  \item Reviewers on a time budget should start with the setup sanity check
        (Installation, Step~2) and the headline \texttt{head8mb} claim, which is
        the cheapest full claim; \texttt{-{}-icount} makes any claim a
        minutes-long end-to-end smoke test.
  \item The full artifact README is at \texttt{experiments/ae/README.md}, and
        the motivation README at \texttt{experiments/ae/motivation/README.md}.
\end{itemize}

%
%  It uses only \section/\subsection, itemize, \texttt, and \url (hyperref).
%  If your template already issues \appendix, delete the \appendix line below.
%
%  FILL-INS (search for <<>>): paper title and, if you mint one, a Zenodo DOI.
% =====================================================================

\appendix
\section{Artifact Appendix}

\subsection{Abstract}

This artifact reproduces the main results of TRAIL, a temporal TLB
prefetcher that makes the page table self-prefetching: each translation
carries its likely successors, so prefetching a translation simultaneously
delivers the next prediction. The artifact evaluates TRAIL against prior
TLB prefetchers (ASP, Stride/NextPage, DP, Recency, ATP, Berti) and the
Perfect-L2TLB upper bound on the Sniper architectural simulator extended
with the MimicOS operating-system model. It reproduces the single-core
head-to-head results at an 8\,MB and a 2\,MB last-level-cache NUCA
(Figures~12 and~13), the in-PTE and side-car payload-budget sweeps
(Tables~5 and~6), the L2-TLB prefetch-queue sensitivity study
(Figure~20), and the four-core multi-programmed study (Figure~22). It also
reproduces the temporal-locality motivation figures (Figures~4, 5, 6, 8,
and~9) directly from page-table-walk dumps. Evaluation is trace-driven, so
no proprietary tools (Intel Pin/SDE) are required: a single script installs
dependencies, builds the simulator, downloads the public traces, and runs a
sanity check; a three-command launch/watch/results workflow then reproduces
each claim and renders the corresponding figure or table.

\subsection{Artifact check-list (meta-information)}

\begin{itemize}
  \item \textbf{Program:} Sniper multi-core simulator with the MimicOS OS
        model and the TRAIL prefetcher (C++17); Python analysis and plotting
        scripts. All included.
  \item \textbf{Compilation:} \texttt{g++} (C++17); the instruction decoder
        (libxed) is bundled and built automatically. No Intel Pin/SDE or
        libtorch is needed for the trace-replay build.
  \item \textbf{Run-time environment:} Linux (tested on Ubuntu 20.04 and
        22.04). Root/\texttt{sudo} only to install system packages.
  \item \textbf{Hardware:} any x86-64 machine. A SLURM cluster is recommended
        for the full run but not required; every step also runs locally.
  \item \textbf{Metrics:} instructions-per-cycle (IPC) and its speedup over a
        no-prefetch baseline. Single-core speedups are geometric means over the
        workload suite; the multi-core metric is an equal-work harmonic mean
        across the four cores.
  \item \textbf{Output:} per-claim markdown tables and rendered figures/tables
        (PDF\,+\,PNG) under \texttt{experiments/ae/ae\_out/}.
  \item \textbf{Experiments:} six simulation claims (\texttt{head8mb},
        \texttt{head2mb}, \texttt{table5}, \texttt{table6}, \texttt{pqsweep},
        \texttt{multicore}) plus the trace-free motivation analysis, driven by
        the provided bash harness.
  \item \textbf{How much disk space required (approximately)?} $\sim$300\,GB
        ($\sim$250\,GB of traces plus the build and results).
  \item \textbf{How much time is needed to prepare workflow (approximately)?}
        $\sim$30--60\,min: a 2--3\,min build plus the (resumable) trace
        download, which dominates and depends on bandwidth.
  \item \textbf{How much time is needed to complete experiments
        (approximately)?} The full set is $\approx$21{,}000 short simulations.
        On a $\sim$1300-core cluster the entire set finishes within
        $\sim$1\,day; the longest single claim (\texttt{table6}) takes
        $\sim$10\,h. A single claim, e.g.\ the headline \texttt{head8mb}, is
        far cheaper; an \texttt{--icount} flag shrinks every job to a
        minutes-long smoke test.
  \item \textbf{Publicly available?} Yes. Source on GitHub; traces and
        page-table-walk dumps as public Hugging Face datasets (no token
        required).
  \item \textbf{Code licenses (if publicly available)?} See the repository
        \texttt{LICENSE}. \emph{<<confirm license, e.g.\ MIT>>}
  \item \textbf{Data licenses (if publicly available)?} Public Hugging Face
        datasets (traces and dumps).
  \item \textbf{Workflow automation framework used?} A self-contained bash
        harness (\texttt{reproduce.sh} for setup, then
        \texttt{ae\_launch.sh}\,/\,\texttt{ae\_watch.sh}\,/\,%
        \texttt{ae\_results.sh}, or \texttt{ae\_run\_all.sh} to run every claim
        in parallel).
  \item \textbf{Archived (provide DOI)?} \emph{<<add Zenodo DOI if minted;
        otherwise the GitHub tag + Hugging Face datasets>>}
\end{itemize}

\subsection{Description}

\subsubsection{How to access}
Clone the artifact branch of the repository:
\begin{quote}\small
\texttt{git clone -{}-branch trail-artifact-release \textbackslash}\\
\texttt{\ \ https://github.com/CMU-SAFARI/Virtuoso.git}
\end{quote}
The traces and page-table-walk dumps are separate public Hugging Face
datasets and are fetched by the setup script:
\begin{itemize}
  \item Traces + trace-lists: \url{https://huggingface.co/datasets/konkanello/trail_traces}
  \item Motivation dumps: \url{https://huggingface.co/datasets/konkanello/trail_ptw_dumps}
\end{itemize}
All artifact material lives under \texttt{experiments/ae/}.

\subsubsection{Hardware dependencies}
An x86-64 machine. The build and the sanity check need $\sim$8\,GB of RAM;
each simulation uses a few GB, so a machine (or cluster) with more memory and
cores runs more jobs in parallel. $\sim$300\,GB of free disk is required for
the traces, build, and results. A SLURM cluster is optional but strongly
recommended to complete the full set in reasonable wall-clock time.

\subsubsection{Software dependencies}
Linux with a C++17 toolchain and Python~3.8+. All system packages, plus the
Hugging Face CLI (trace download) and matplotlib (figures), are installed by
\texttt{experiments/ae/lib/install\_deps.sh} on Debian/Ubuntu (equivalents
exist on other distributions). Network access is needed to download the traces
and to fetch the decoder during the build. No license-gated or proprietary
tools are required.

\subsubsection{Benchmarks / data sets}
The workload suite spans server, graph-analytics, machine-learning, HPC, and
industrial traces. The public trace dataset ships the trace files together
with the trace-lists; the setup script wires them into the repository. The
motivation figures use a separate $\sim$3\,GB dataset of per-workload
page-table-walk dumps.

\subsection{Installation}

\begin{enumerate}
  \item Install dependencies:
\begin{quote}\small
\texttt{bash experiments/ae/lib/install\_deps.sh}
\end{quote}
  \item Build, download traces, and sanity-check (resumable; re-run to resume
        after any interruption):
\begin{quote}\small
\texttt{bash experiments/ae/reproduce.sh -{}-skip-deps}
\end{quote}
        This performs a trace-replay build of the simulator
        (\texttt{make SNIPER\_TRACE\_ONLY=1 replay}; no Pin/SDE/libtorch),
        downloads the public trace dataset, resolves the trace-lists into
        \texttt{experiments/vm\_tlist/}, and runs a few short simulations on
        randomly chosen traces, printing \texttt{[PASS] <trace> IPC=...} for
        each. It then stops and prints the commands for Step~3. This step alone
        confirms that the artifact builds and runs end-to-end.
\end{enumerate}

\subsection{Experiment workflow}

Each claim maps to one paper figure or table. Results and rendered figures are
written to \texttt{experiments/ae/ae\_out/}.

\emph{All claims in parallel (recommended on a cluster):}
\begin{quote}\small
\texttt{bash experiments/ae/ae\_run\_all.sh -{}-mode slurm \textbackslash}\\
\texttt{\ \ [-{}-partitions <p>]}\qquad\textit{\# launch every claim}\\
\texttt{bash experiments/ae/ae\_run\_all.sh -{}-status}\quad\textit{\# progress}\\
\texttt{bash experiments/ae/ae\_run\_all.sh -{}-results}\quad\textit{\# figures/tables}
\end{quote}
\texttt{-{}-partitions} is optional (omit it to use the cluster default; nothing
is then passed to \texttt{sbatch}).

\emph{One claim at a time (launch $\rightarrow$ watch $\rightarrow$ results):}
\begin{quote}\small
\texttt{bash experiments/ae/ae\_launch.sh\ -{}-claim head8mb -{}-mode slurm}\\
\texttt{bash experiments/ae/ae\_watch.sh\ \ -{}-claim head8mb}\\
\texttt{bash experiments/ae/ae\_results.sh -{}-claim head8mb}\quad\textit{\# after .DONE}
\end{quote}
The watcher self-detaches and maintains a human-readable
\texttt{ae\_out/<claim>.status}; when every job is accounted for it writes a
pass/fail report to \texttt{ae\_out/<claim>.DONE}, and \texttt{ae\_results.sh}
then parses the results and renders the figure/table. Launching is
validity-based and resumable: only jobs without a valid \texttt{sim.stats} are
(re)submitted, so an interrupted run resumes cleanly. On a single machine, use
\texttt{-{}-mode local} one claim at a time; \texttt{-{}-icount N} shrinks every
job's instruction budget for a quick end-to-end test.

The claim-to-paper mapping is:

\begin{center}\small
\begin{tabular}{@{}llr@{}}
\toprule
\texttt{-{}-claim} & reproduces & jobs \\
\midrule
\texttt{head8mb}   & Fig.~12 (bottom, 8\,MB) + Fig.~13 & 2761 \\
\texttt{head2mb}   & Fig.~12 (top, 2\,MB)              & 2761 \\
\texttt{table5}    & Table~5 (in-PTE payload sweep)    & 5271 \\
\texttt{table6}    & Table~6 (side-car payload sweep)  & 6275 \\
\texttt{pqsweep}   & Fig.~20 (PQ-size sensitivity)     & 3012 \\
\texttt{multicore} & Fig.~22 (4-core, 100 mixes)       & 1000 \\
\midrule
\textbf{all}       &                                   & \textbf{21080} \\
\bottomrule
\end{tabular}
\end{center}

\subsection{Evaluation and expected results}

Each claim writes a markdown table (\texttt{ae\_out/<claim>.md}) alongside the
rendered figure/table. Exact numbers vary slightly with the machine, but every
claim should closely match the paper, and in all cases \textbf{TRAIL should
beat every prior prefetcher} and move toward the Perfect-L2TLB upper bound.
Speedups are geometric means over the suite (multi-core is an equal-work
harmonic mean across the four cores):

\begin{center}\small
\begin{tabular}{@{}ll@{}}
\toprule
claim & expected (approx.) \\
\midrule
\texttt{head8mb} (Fig.~12, 8\,MB) & TRAIL $\approx +4.7\%$ over no-prefetch \\
                                  & ($+2.4\%$ over ASP); Perfect $\approx +11\%$ \\
\texttt{head2mb} (Fig.~12, 2\,MB) & TRAIL $\approx +5.1\%$ ($+2.8\%$ over ASP); \\
                                  & Perfect $\approx +16.5\%$ \\
\texttt{table5}  (Table~5)        & in-PTE TRAIL grows with the budget, \\
                                  & up to $\approx +2.4\%$ over ASP \\
\texttt{table6}  (Table~6)        & side-car TRAIL grows with the budget, \\
                                  & up to $\approx +2.5\%$ over ASP \\
\texttt{pqsweep} (Fig.~20)        & TRAIL $\approx +2.2$--$2.4\%$ over same- \\
                                  & size ASP across every PQ size (64--1024) \\
\texttt{multicore} (Fig.~22)      & TRAIL $\approx +11\%$ harmonic mean (best \\
                                  & prior $\approx +5\%$); Perfect $\approx +23\%$ \\
\bottomrule
\end{tabular}
\end{center}

\subsection{Motivation figures (Figures 4, 5, 6, 8, 9)}

The temporal-locality motivation figures are produced from the page-table-walk
dumps, without running any simulation. Download the dumps and run the analysis
(SLURM or local, using up to \texttt{\$(nproc)} cores):
\begin{quote}\small
\texttt{hf download konkanello/trail\_ptw\_dumps \textbackslash}\\
\texttt{\ \ -{}-repo-type dataset -{}-local-dir ./ptw\_bundle}\\[2pt]
\texttt{bash experiments/ae/motivation/run\_motivation.sh \textbackslash}\\
\texttt{\ \ -{}-dumps ./ptw\_bundle -{}-mode local -{}-jobs \$(nproc)}
\end{quote}
This writes \texttt{figure4}--\texttt{figure9} (PDF\,+\,PNG) under
\texttt{experiments/ae/motivation/motivation\_out/}. The analysis is resumable:
one JSON per workload, re-plotted on any re-run.

\subsection{Notes}

\begin{itemize}
  \item Every phase is resumable. Re-running \texttt{reproduce.sh}, a launch, or
        the motivation driver skips whatever is already complete.
  \item Reviewers on a time budget should start with the setup sanity check
        (Installation, Step~2) and the headline \texttt{head8mb} claim, which is
        the cheapest full claim; \texttt{-{}-icount} makes any claim a
        minutes-long end-to-end smoke test.
  \item The full artifact README is at \texttt{experiments/ae/README.md}, and
        the motivation README at \texttt{experiments/ae/motivation/README.md}.
\end{itemize}

%
%  It uses only \section/\subsection, itemize, \texttt, and \url (hyperref).
%  If your template already issues \appendix, delete the \appendix line below.
%
%  FILL-INS (search for <<>>): paper title and, if you mint one, a Zenodo DOI.
% =====================================================================

\appendix
\section{Artifact Appendix}

\subsection{Abstract}

This artifact reproduces the main results of TRAIL, a temporal TLB
prefetcher that makes the page table self-prefetching: each translation
carries its likely successors, so prefetching a translation simultaneously
delivers the next prediction. The artifact evaluates TRAIL against prior
TLB prefetchers (ASP, Stride/NextPage, DP, Recency, ATP, Berti) and the
Perfect-L2TLB upper bound on the Sniper architectural simulator extended
with the MimicOS operating-system model. It reproduces the single-core
head-to-head results at an 8\,MB and a 2\,MB last-level-cache NUCA
(Figures~12 and~13), the in-PTE and side-car payload-budget sweeps
(Tables~5 and~6), the L2-TLB prefetch-queue sensitivity study
(Figure~20), and the four-core multi-programmed study (Figure~22). It also
reproduces the temporal-locality motivation figures (Figures~4, 5, 6, 8,
and~9) directly from page-table-walk dumps. Evaluation is trace-driven, so
no proprietary tools (Intel Pin/SDE) are required: a single script installs
dependencies, builds the simulator, downloads the public traces, and runs a
sanity check; a three-command launch/watch/results workflow then reproduces
each claim and renders the corresponding figure or table.

\subsection{Artifact check-list (meta-information)}

\begin{itemize}
  \item \textbf{Program:} Sniper multi-core simulator with the MimicOS OS
        model and the TRAIL prefetcher (C++17); Python analysis and plotting
        scripts. All included.
  \item \textbf{Compilation:} \texttt{g++} (C++17); the instruction decoder
        (libxed) is bundled and built automatically. No Intel Pin/SDE or
        libtorch is needed for the trace-replay build.
  \item \textbf{Run-time environment:} Linux (tested on Ubuntu 20.04 and
        22.04). Root/\texttt{sudo} only to install system packages.
  \item \textbf{Hardware:} any x86-64 machine. A SLURM cluster is recommended
        for the full run but not required; every step also runs locally.
  \item \textbf{Metrics:} instructions-per-cycle (IPC) and its speedup over a
        no-prefetch baseline. Single-core speedups are geometric means over the
        workload suite; the multi-core metric is an equal-work harmonic mean
        across the four cores.
  \item \textbf{Output:} per-claim markdown tables and rendered figures/tables
        (PDF\,+\,PNG) under \texttt{experiments/ae/ae\_out/}.
  \item \textbf{Experiments:} six simulation claims (\texttt{head8mb},
        \texttt{head2mb}, \texttt{table5}, \texttt{table6}, \texttt{pqsweep},
        \texttt{multicore}) plus the trace-free motivation analysis, driven by
        the provided bash harness.
  \item \textbf{How much disk space required (approximately)?} $\sim$300\,GB
        ($\sim$250\,GB of traces plus the build and results).
  \item \textbf{How much time is needed to prepare workflow (approximately)?}
        $\sim$30--60\,min: a 2--3\,min build plus the (resumable) trace
        download, which dominates and depends on bandwidth.
  \item \textbf{How much time is needed to complete experiments
        (approximately)?} The full set is $\approx$21{,}000 short simulations.
        On a $\sim$1300-core cluster the entire set finishes within
        $\sim$1\,day; the longest single claim (\texttt{table6}) takes
        $\sim$10\,h. A single claim, e.g.\ the headline \texttt{head8mb}, is
        far cheaper; an \texttt{--icount} flag shrinks every job to a
        minutes-long smoke test.
  \item \textbf{Publicly available?} Yes. Source on GitHub; traces and
        page-table-walk dumps as public Hugging Face datasets (no token
        required).
  \item \textbf{Code licenses (if publicly available)?} See the repository
        \texttt{LICENSE}. \emph{<<confirm license, e.g.\ MIT>>}
  \item \textbf{Data licenses (if publicly available)?} Public Hugging Face
        datasets (traces and dumps).
  \item \textbf{Workflow automation framework used?} A self-contained bash
        harness (\texttt{reproduce.sh} for setup, then
        \texttt{ae\_launch.sh}\,/\,\texttt{ae\_watch.sh}\,/\,%
        \texttt{ae\_results.sh}, or \texttt{ae\_run\_all.sh} to run every claim
        in parallel).
  \item \textbf{Archived (provide DOI)?} \emph{<<add Zenodo DOI if minted;
        otherwise the GitHub tag + Hugging Face datasets>>}
\end{itemize}

\subsection{Description}

\subsubsection{How to access}
Clone the artifact branch of the repository:
\begin{quote}\small
\texttt{git clone -{}-branch trail-artifact-release \textbackslash}\\
\texttt{\ \ https://github.com/CMU-SAFARI/Virtuoso.git}
\end{quote}
The traces and page-table-walk dumps are separate public Hugging Face
datasets and are fetched by the setup script:
\begin{itemize}
  \item Traces + trace-lists: \url{https://huggingface.co/datasets/konkanello/trail_traces}
  \item Motivation dumps: \url{https://huggingface.co/datasets/konkanello/trail_ptw_dumps}
\end{itemize}
All artifact material lives under \texttt{experiments/ae/}.

\subsubsection{Hardware dependencies}
An x86-64 machine. The build and the sanity check need $\sim$8\,GB of RAM;
each simulation uses a few GB, so a machine (or cluster) with more memory and
cores runs more jobs in parallel. $\sim$300\,GB of free disk is required for
the traces, build, and results. A SLURM cluster is optional but strongly
recommended to complete the full set in reasonable wall-clock time.

\subsubsection{Software dependencies}
Linux with a C++17 toolchain and Python~3.8+. All system packages, plus the
Hugging Face CLI (trace download) and matplotlib (figures), are installed by
\texttt{experiments/ae/lib/install\_deps.sh} on Debian/Ubuntu (equivalents
exist on other distributions). Network access is needed to download the traces
and to fetch the decoder during the build. No license-gated or proprietary
tools are required.

\subsubsection{Benchmarks / data sets}
The workload suite spans server, graph-analytics, machine-learning, HPC, and
industrial traces. The public trace dataset ships the trace files together
with the trace-lists; the setup script wires them into the repository. The
motivation figures use a separate $\sim$3\,GB dataset of per-workload
page-table-walk dumps.

\subsection{Installation}

\begin{enumerate}
  \item Install dependencies:
\begin{quote}\small
\texttt{bash experiments/ae/lib/install\_deps.sh}
\end{quote}
  \item Build, download traces, and sanity-check (resumable; re-run to resume
        after any interruption):
\begin{quote}\small
\texttt{bash experiments/ae/reproduce.sh -{}-skip-deps}
\end{quote}
        This performs a trace-replay build of the simulator
        (\texttt{make SNIPER\_TRACE\_ONLY=1 replay}; no Pin/SDE/libtorch),
        downloads the public trace dataset, resolves the trace-lists into
        \texttt{experiments/vm\_tlist/}, and runs a few short simulations on
        randomly chosen traces, printing \texttt{[PASS] <trace> IPC=...} for
        each. It then stops and prints the commands for Step~3. This step alone
        confirms that the artifact builds and runs end-to-end.
\end{enumerate}

\subsection{Experiment workflow}

Each claim maps to one paper figure or table. Results and rendered figures are
written to \texttt{experiments/ae/ae\_out/}.

\emph{All claims in parallel (recommended on a cluster):}
\begin{quote}\small
\texttt{bash experiments/ae/ae\_run\_all.sh -{}-mode slurm \textbackslash}\\
\texttt{\ \ [-{}-partitions <p>]}\qquad\textit{\# launch every claim}\\
\texttt{bash experiments/ae/ae\_run\_all.sh -{}-status}\quad\textit{\# progress}\\
\texttt{bash experiments/ae/ae\_run\_all.sh -{}-results}\quad\textit{\# figures/tables}
\end{quote}
\texttt{-{}-partitions} is optional (omit it to use the cluster default; nothing
is then passed to \texttt{sbatch}).

\emph{One claim at a time (launch $\rightarrow$ watch $\rightarrow$ results):}
\begin{quote}\small
\texttt{bash experiments/ae/ae\_launch.sh\ -{}-claim head8mb -{}-mode slurm}\\
\texttt{bash experiments/ae/ae\_watch.sh\ \ -{}-claim head8mb}\\
\texttt{bash experiments/ae/ae\_results.sh -{}-claim head8mb}\quad\textit{\# after .DONE}
\end{quote}
The watcher self-detaches and maintains a human-readable
\texttt{ae\_out/<claim>.status}; when every job is accounted for it writes a
pass/fail report to \texttt{ae\_out/<claim>.DONE}, and \texttt{ae\_results.sh}
then parses the results and renders the figure/table. Launching is
validity-based and resumable: only jobs without a valid \texttt{sim.stats} are
(re)submitted, so an interrupted run resumes cleanly. On a single machine, use
\texttt{-{}-mode local} one claim at a time; \texttt{-{}-icount N} shrinks every
job's instruction budget for a quick end-to-end test.

The claim-to-paper mapping is:

\begin{center}\small
\begin{tabular}{@{}llr@{}}
\toprule
\texttt{-{}-claim} & reproduces & jobs \\
\midrule
\texttt{head8mb}   & Fig.~12 (bottom, 8\,MB) + Fig.~13 & 2761 \\
\texttt{head2mb}   & Fig.~12 (top, 2\,MB)              & 2761 \\
\texttt{table5}    & Table~5 (in-PTE payload sweep)    & 5271 \\
\texttt{table6}    & Table~6 (side-car payload sweep)  & 6275 \\
\texttt{pqsweep}   & Fig.~20 (PQ-size sensitivity)     & 3012 \\
\texttt{multicore} & Fig.~22 (4-core, 100 mixes)       & 1000 \\
\midrule
\textbf{all}       &                                   & \textbf{21080} \\
\bottomrule
\end{tabular}
\end{center}

\subsection{Evaluation and expected results}

Each claim writes a markdown table (\texttt{ae\_out/<claim>.md}) alongside the
rendered figure/table. Exact numbers vary slightly with the machine, but every
claim should closely match the paper, and in all cases \textbf{TRAIL should
beat every prior prefetcher} and move toward the Perfect-L2TLB upper bound.
Speedups are geometric means over the suite (multi-core is an equal-work
harmonic mean across the four cores):

\begin{center}\small
\begin{tabular}{@{}ll@{}}
\toprule
claim & expected (approx.) \\
\midrule
\texttt{head8mb} (Fig.~12, 8\,MB) & TRAIL $\approx +4.7\%$ over no-prefetch \\
                                  & ($+2.4\%$ over ASP); Perfect $\approx +11\%$ \\
\texttt{head2mb} (Fig.~12, 2\,MB) & TRAIL $\approx +5.1\%$ ($+2.8\%$ over ASP); \\
                                  & Perfect $\approx +16.5\%$ \\
\texttt{table5}  (Table~5)        & in-PTE TRAIL grows with the budget, \\
                                  & up to $\approx +2.4\%$ over ASP \\
\texttt{table6}  (Table~6)        & side-car TRAIL grows with the budget, \\
                                  & up to $\approx +2.5\%$ over ASP \\
\texttt{pqsweep} (Fig.~20)        & TRAIL $\approx +2.2$--$2.4\%$ over same- \\
                                  & size ASP across every PQ size (64--1024) \\
\texttt{multicore} (Fig.~22)      & TRAIL $\approx +11\%$ harmonic mean (best \\
                                  & prior $\approx +5\%$); Perfect $\approx +23\%$ \\
\bottomrule
\end{tabular}
\end{center}

\subsection{Motivation figures (Figures 4, 5, 6, 8, 9)}

The temporal-locality motivation figures are produced from the page-table-walk
dumps, without running any simulation. Download the dumps and run the analysis
(SLURM or local, using up to \texttt{\$(nproc)} cores):
\begin{quote}\small
\texttt{hf download konkanello/trail\_ptw\_dumps \textbackslash}\\
\texttt{\ \ -{}-repo-type dataset -{}-local-dir ./ptw\_bundle}\\[2pt]
\texttt{bash experiments/ae/motivation/run\_motivation.sh \textbackslash}\\
\texttt{\ \ -{}-dumps ./ptw\_bundle -{}-mode local -{}-jobs \$(nproc)}
\end{quote}
This writes \texttt{figure4}--\texttt{figure9} (PDF\,+\,PNG) under
\texttt{experiments/ae/motivation/motivation\_out/}. The analysis is resumable:
one JSON per workload, re-plotted on any re-run.

\subsection{Notes}

\begin{itemize}
  \item Every phase is resumable. Re-running \texttt{reproduce.sh}, a launch, or
        the motivation driver skips whatever is already complete.
  \item Reviewers on a time budget should start with the setup sanity check
        (Installation, Step~2) and the headline \texttt{head8mb} claim, which is
        the cheapest full claim; \texttt{-{}-icount} makes any claim a
        minutes-long end-to-end smoke test.
  \item The full artifact README is at \texttt{experiments/ae/README.md}, and
        the motivation README at \texttt{experiments/ae/motivation/README.md}.
\end{itemize}
